% =====================================================================
\parte{Parte A. Cuestiones}{cap:test}
% =====================================================================

Veinte preguntas de opción múltiple, de 1,5 puntos cada una (30 puntos en total). Cada pregunta tiene una sola respuesta correcta; las respuestas erróneas no restan. Tiempo orientativo: 30 minutos.

\subsection*{Módulo 1. Modelos de lenguaje probabilísticos}

\pregunta{La regla de la cadena $P(x_1,\dots,x_T)=\prod_{t=1}^{T}P(x_t\mid x_1,\dots,x_{t-1})$:}
{es una aproximación que solo es válida para modelos de bigramas.}
{es una identidad exacta para cualquier distribución sobre secuencias.}
{exige que los tokens sean independientes entre sí.}
{solo es válida si todas las probabilidades son estrictamente positivas.}

\pregunta{Con un vocabulario de 28 tokens (27 letras y el símbolo \fin), ¿cuántas casillas tiene la tabla de conteos de un modelo de trigramas?}
{$784$}
{$2187$}
{$21\,952$}
{$614\,656$}

\pregunta{Un modelo asigna probabilidad $1/4$ a cada token de un texto de validación. Su perplejidad sobre ese texto es:}
{$0{,}25$}
{$2$}
{$4$}
{$\ln 4\approx 1{,}386$}

\pregunta{Cuando la temperatura tiende a infinito, $T\to\infty$, la distribución $\softmax_T(z)$ de un vector de logits $z\in\R^{K}$ tiende a:}
{un vector one-hot en la posición del logit máximo.}
{la distribución uniforme $(1/K,\dots,1/K)$.}
{$\softmax(z)$, la distribución a temperatura $T=1$.}
{el vector nulo.}

\subsection*{Módulo 2. Álgebra lineal I}

\pregunta{La similitud coseno entre $u=(1,1)$ y $v=(3,0)$ vale:}
{$0$}
{$0{,}5$}
{$0{,}707$}
{$1$}

\pregunta{Sea $E\in\R^{K\times d}$ la matriz de embeddings, con una fila por token, y sea $e_i\in\R^{K}$ el vector one-hot del token $i$. El producto $e_i\T E$ es:}
{la columna $i$ de $E$.}
{la fila $i$ de $E$.}
{la suma de todas las filas de $E$.}
{un escalar.}

\subsection*{Módulo 3. Álgebra lineal II}

\pregunta{En NumPy, un tensor de forma $(8,128,512)$ multiplicado mediante el operador \texttt{@} por una matriz de forma $(512,\,50\,000)$ produce un tensor de forma:}
{$(8,\,128,\,50\,000)$}
{$(8,\,512,\,50\,000)$}
{$(128,\,50\,000)$}
{ninguna: se produce un error de dimensiones.}

\pregunta{La matriz $\begin{pmatrix}2&1\\4&2\end{pmatrix}$:}
{tiene rango~2 y es invertible.}
{tiene rango~1 y uno de sus autovalores es~0.}
{tiene rango~0.}
{tiene determinante~8.}

\pregunta{Una matriz tiene valores singulares $5$, $3$ y $1$. El error, en norma de Frobenius, de su mejor aproximación de rango~1 es:}
{$1$}
{$3$}
{$\sqrt{10}\approx 3{,}16$}
{$4$}

\pregunta{Antes de aplicar el análisis de componentes principales (PCA) se centran los datos porque:}
{sin centrar, la SVD de la matriz de datos no existe.}
{sin centrar, la primera componente tiende a apuntar hacia la media de los datos en lugar de hacia la dirección de máxima varianza.}
{al centrar, los autovalores de la covarianza son negativos.}
{al centrar, la matriz de covarianza es diagonal.}

\pregunta{Se suman dos arrays con las reglas de \emph{broadcasting} de NumPy. ¿Cuál de estas combinaciones de formas produce un error?}
{$(16,128,512)$ y $(512,)$}
{$(16,128,512)$ y $(128,1)$}
{$(16,128,512)$ y $(16,128)$}
{$(16,128,512)$ y $(1,1,512)$}

\subsection*{Módulo 4. Cálculo para el aprendizaje automático}

\pregunta{El gradiente de $f(x,y)=x^{2}y$ en el punto $(1,3)$ es:}
{$(2,1)$}
{$(6,1)$}
{$(6,3)$}
{$(3,2)$}

\pregunta{Sean $p=\softmax(z)$ e $y$ el vector one-hot de la clase correcta $c$. Para la pérdida $L=-\ln\softmax(z)_c$, el gradiente $\nabla_z L$ es:}
{$p$}
{$y-p$}
{$p-y$}
{$-\ln p$}

\pregunta{Para una función $f$ de clase $C^{3}$, el error de truncamiento de la diferencia centrada $[f(x+h)-f(x-h)]/(2h)$ como aproximación de $f'(x)$ es del orden de:}
{$h$}
{$h^{2}$}
{$1/h$}
{no depende de $h$.}

\subsection*{Módulo 5. Probabilidad y estadística}

\pregunta{Una variable de Bernoulli de parámetro $p=0{,}2$ tiene varianza:}
{$0{,}04$}
{$0{,}16$}
{$0{,}2$}
{$0{,}8$}

\pregunta{El estimador de máxima verosimilitud de los parámetros de una distribución categórica a partir de los conteos $N_1,\dots,N_K$ es:}
{$\hat p_j=1/K$}
{$\hat p_j=(N_j+1)/\big(\sum_k N_k+K\big)$}
{$\hat p_j=N_j/\sum_k N_k$}
{$\hat p_j=\sqrt{N_j}/\sum_k\sqrt{N_k}$}

\pregunta{Si se multiplica por~4 el tamaño de un minilote de ejemplos \emph{i.i.d.}, la desviación típica de la pérdida media del minilote:}
{se divide por~4.}
{se divide por~2.}
{no cambia.}
{se multiplica por~2.}

\subsection*{Módulo 6. Teoría de la información}

\pregunta{La entropía de una distribución uniforme sobre 8 símbolos es:}
{1 bit}
{2 bits}
{3 bits}
{8 bits}

\pregunta{Sobre la divergencia de Kullback--Leibler $\KL{p}{q}$ entre distribuciones discretas es cierto que:}
{es simétrica: $\KL{p}{q}=\KL{q}{p}$.}
{puede ser negativa si $q$ se parece mucho a $p$.}
{es mayor o igual que~0 y vale~0 si y solo si $p=q$.}
{es igual a la entropía de $q$.}

\pregunta{Un texto tiene entropía $H(p)=2$ bits por carácter y un modelo $q$ alcanza sobre él una entropía cruzada $H(p,q)=2{,}5$ bits por carácter. La divergencia $\KL{p}{q}$ vale:}
{$0{,}5$ bits}
{$1{,}25$ bits}
{$4{,}5$ bits}
{no puede determinarse con estos datos.}
